\section*{Chapter \ref{ch:b3:facies-models} --- Facies Models}

\begin{solution}{exo:b3:facies-models:1}
From the sea to the land: offshore, offshore transition with storm sands,
\omterm{def:b3:facies-models:shoreface}{shoreface}, beach, a coastal plain with swamps. Each environment came to lie
over its seaward neighbour: the shoreline moved seaward, a regression. A
transgression would have stacked them the other way, coastal plain at the
base and offshore mud at the top, often with an erosion surface where the
\omterm{def:b3:facies-models:shoreface}{shoreface} passed.
\end{solution}

\begin{solution}{exo:b3:facies-models:2}
(a) A channel fill or \omterm{def:b3:facies-models:rivers}{point bar}: a pebbly scour at the base, cross-beds,
then ripples and rooted mud. (b) A prograding \omterm{def:b3:facies-models:shoreface}{shoreface} or delta front:
mud with storm sand beds and \omterm{def:b3:facies-models:hcs}{hummocky cross-stratification} below,
cross-bedded \omterm{def:b3:facies-models:shoreface}{shoreface} sand and beach above. (c) Stacked braided-river
sand, an aeolian dune or a carbonate shoal: cross-bed sets of constant grain
size, with large dune foresets for the aeolian case or skeletal and oolitic
grains for the shoal.
\end{solution}

\begin{solution}{exo:b3:facies-models:3}
$\phi = (2.65 - 2.40)/(2.65 - 1.0) = 0.152$; with brine,
$0.25/1.55 = 0.161$.
\end{solution}

\begin{solution}{exo:b3:facies-models:4}
$e^{-2\pi z/\lambda} = e^{-\pi} = 0.043$ at $z = \lambda/2$: about
\qty{38}{m} and \qty{55}{m}. At \qty{20}{m}: $e^{-2\pi\times20/75} = 0.19$
and $e^{-2\pi\times20/110} = 0.32$ of the surface motion.
\end{solution}

\begin{solution}{exo:b3:facies-models:5}
$(60 - 47.6)/(189 - 47.6) = 0.088$; $(60 - 52.6)/(189 - 52.6) = 0.054$. A
matrix value uncertain by \qty{5}{\micro\second/ft} changes the porosity by
a third: the time-average equation needs a matrix value checked on cores.
\end{solution}

\begin{solution}{exo:b3:facies-models:6}
$(6.0 + 3.5 + 8.0 + 2.5)/40 = 20/40 = 0.50$. Pore volume: $20\times0.22
\times10^6 = \qty{4.4e6}{m^3}$ per square kilometre.
\end{solution}

\begin{solution}{exo:b3:facies-models:7}
Near the canyon: Ta only, the upper divisions removed by the next flow; down
the fan: Tc (ripples) and Td (silt) under Te (mud). The current loses speed
and its coarse load with distance; only the fine fraction travels
\qty{30}{km}, and it is deposited by a slower flow that builds ripples, not
by fallout from a dense suspension.
\end{solution}

\begin{solution}{exo:b3:facies-models:8}
A \omterm{def:b3:facies-models:rivers}{point bar} of a \omterm{def:b3:facies-models:rivers}{meandering river}: lag on the scoured channel floor,
cross-bedded sand of the deep channel, rippled sand of the top of the bar,
floodplain mud with roots. The succession records the channel from floor to
bank, so the channel was about \qty{9}{m} deep (a little less, after
compaction of the mud). Not across the base: it is an erosion surface, and
what lay below may have been any older deposit.
\end{solution}

\begin{solution}{exo:b3:facies-models:9}
$60/40\,000 = 0.0015$, \ang{0.09}: a ramp. $300/1500 = 0.20$, \ang{11.3},
after a flat top \qty{5}{m} deep: a \omterm{def:b3:facies-models:ramp}{rimmed shelf}, with reef or shoal grains
at the margin and blocks and grain flows on the slope below.
\end{solution}

\begin{solution}{exo:b3:facies-models:10}
The belt takes a time $w/v$ to pass a fixed place, during which the floor
rises relative to sea level at $A$ under it: the sand fills $h = A\,w/v$.
$2000\times0.2/20 = \qty{20}{m}$. With $A = 0$ the \omterm{def:b3:facies-models:shoreface}{shoreface} leaves only
the thickness of its own profile, a thin sheet, as the belts slide over one
another without aggrading. The ratio $A/v$ is the tangent of the trajectory
angle of the proposition of Chapter~\ref{ch:b3:sequence-stratigraphy}: the
thickness of each facies is its width times that tangent.
\end{solution}

\begin{solution}{exo:b3:facies-models:11}
Density: $(2.65 - 2.30)/1.65 = 0.212$; sonic: $(86 - 52.6)/136.4 = 0.245$.
With $\phi = 0.245$, $\rho_{\mathrm{ma}} = (2.30 - 0.245)/(1 - 0.245) =
\qty{2.72}{g/cm^3}$: denser than quartz, as for a sandstone cemented by
calcite (2.71) or holding calcite grains.
\end{solution}

\begin{solution}{exo:b3:facies-models:12}
Offshore mud on a coal brings non-adjacent environments into contact: a
\omterm{def:b3:sequence-stratigraphy:parasequence}{flooding surface}, here a \omterm{def:b3:sequence-stratigraphy:surfaces}{transgressive surface} of wave erosion, the shell
debris a lag. As the shoreline moved landward, the waves cut the floor down
to about the \omterm{def:b3:facies-models:shoreface}{fair-weather wave base}: the beach, the upper \omterm{def:b3:facies-models:shoreface}{shoreface} and the
top of the coastal plain, up to about \qty{10}{m} of section, may be
missing.
\end{solution}

\begin{solution}{pb:b3:facies-models:1}
\textbf{1.} Chalk (cold, low density), shale from \qty{766}{m} (hot), limestone from
\qty{866}{m} (cold, dense, fast), sandstone from \qty{1046}{m} (bell
shapes), schist from \qty{1166}{m}.

\textbf{2.} Chalk \qty{766}{m}, shale 100, limestone 180, sandstone
\qty{120}{m}; base of the cover at $441.0 - 1166 = \qty{-725}{m}$.

\textbf{3.} $(2.71 - 2.03)/1.71 = 0.40$ near the surface, $(2.71 -
2.36)/1.71 = 0.20$ at \qty{750}{m}: the chalk loses half its porosity
under the weight of the overburden, by compaction.

\textbf{4.} Its grains are clays of uncertain density, and its water is
partly bound to them: the equation's assumptions fail.

\textbf{5.} Crystalline basement: the cover is crossed, and no reservoir is
expected below.

\textbf{6.} $(2.71 - 2.573)/1.71 = 0.080$.

\textbf{7.} $(58.9 - 47.6)/141.4 = 0.080$: the same.

\textbf{8.} $(2.65 - 2.32)/1.65 = 0.200$.

\textbf{9.} $(79.9 - 52.6)/136.4 = 0.200$: the same.

\textbf{10.} $0.33/1.55 = 0.213$, one point and a third more.

\textbf{11.} The sandstone, with \qty{20}{\percent} porosity in its sands;
the shale above the limestone seals it, and the floodplain muds divide the
sandstone into compartments.

\textbf{12.} A bell: low gamma ray in the cross-bedded sand, rising through
the fining interval to the high values of the rooted mud, on a sharp base.

\textbf{13.} A \omterm{def:b3:facies-models:association}{fining-upward succession}: channel-floor lag on an erosional
base, channel dunes (cross-beds), top of the \omterm{def:b3:facies-models:rivers}{point bar} (ripples), floodplain
(rooted mud).

\textbf{14.} The channel, its \omterm{def:b3:facies-models:rivers}{point bar} on the inner bank, levées and a
floodplain with soils, side by side.

\textbf{15.} About the thickness of sand and fining interval, $10.4 + 3.2 =
\qty{13.6}{m}$, a little more before compaction.

\textbf{16.} A \omterm{def:b3:facies-models:rivers}{meandering river} whose channel belt moved repeatedly by
\omterm{def:b3:facies-models:rivers}{avulsion} while the basin subsided, each new channel cutting into the
floodplain left by the last.

\textbf{17.} No: each base is erosional; the floodplain mud below it and the
channel sand above it need not have been neighbours.

\textbf{18.} $(60 - 30)/(105 - 30) = 0.40$.

\textbf{19.} Lowest: $14\times(0.60 + 0.40\times0.25) = \qty{9.80}{m}$;
thickest: $16\times(0.65 + 0.40\times0.20) = \qty{11.68}{m}$.

\textbf{20.} Clean sand \qty{58.8}{m}; fining intervals \qty{33.4}{m}; mud
\qty{27.8}{m}.

\textbf{21.} Net $= 58.8 + 0.40\times33.4 = \qty{72.2}{m}$; net-to-gross
$72.2/120 = 0.60$.

\textbf{22.} Fractions 0.20 and 0.60: \qty{65.5}{m} and \qty{78.8}{m}, a
spread of \qty{13}{m} from the choice of cutoff alone.

\textbf{23.} $72.2\times0.20\times10^6 = \qty{1.44e7}{m^3}$ per square
kilometre.

\textbf{24.} The net sand thickness of the Triassic sandstone at Varden-2 is
\textbf{\qty[text-series-to-math=true]{72.2}{m}}.
\end{solution}
